\section{Discussion}
\label{sec:discussion}

\subsection{Workload-Dependent Behavior}

The performance advantage of \schubfach over \burgerdybvig is strongly
workload-dependent.  On number-heavy inputs (e.g., \texttt{numeric-boundary},
\texttt{array-2048}), \schubfach achieves 1.4--2.1$\times$ speedup due
to zero heap allocation in the digit-generation path.  On string-dominant
inputs (e.g., \texttt{long-string}, \texttt{unicode}), the advantage
narrows to $<$5\% because formatting cost is dominated by string
processing and key sorting, amortizing the digit-generation difference.

\subsection{Allocation Cost Attribution}

The \burgerdybvig implementation's \texttt{sync.Pool} partially mitigates
allocation cost by reusing big-integer buffers.  Without pooling, the
difference would be larger.  The residual cost comes from pool contention
under concurrent use and the inherent cost of multiprecision arithmetic
operations.

\subsection{Profile-Level Analysis}

CPU profiling confirms that digit generation consumes 15--40\% of total
canonicalization time depending on workload number density.  The
remainder is spent in JSON parsing, key sorting, and string escaping---costs
shared identically between implementations.

\subsection{Practical Implications}

For applications where canonicalization throughput is critical (e.g.,
high-frequency signature verification of number-heavy payloads),
\schubfach offers meaningful improvement.  For typical mixed-content JSON,
the choice of digit-generation algorithm has minimal practical impact on
end-to-end latency.
